%!TEX root = ../dissertation.tex

\chapter{Introduction}\label{chapter intro}
\section{A Brief History}\label{section history}
In 1851 in Victorian England, Francis Guthrie shared a simple observation with his younger brother, challenging mathematicians for decades and laying the foundations for a new subfield in graph theory. As the story goes, Guthrie, an English lawyer, botanist, and mathematician, was coloring maps of England when he noticed the ``greatest necessary number of colours to be used in colouring a map so as to avoid identity of colour in lineally contiguous districts is four''~\cite{FGbro}. For a given map, if we associate each region with a vertex and include an edge between two vertices exactly when the two corresponding regions are adjacent, then we have constructed a graph. A \textit{graph} $G = (V, E)$ refers to a set of vertices $V(G)$ and a set of edges $E(G)$ between these vertices. A few examples of graphs are displayed in Figure~\ref{fig:graphs}.
\begin{figure}
    \centering
    \begin{tikzpicture}[scale=1]

% horizontal spacing
\def\dx{4}

% ================= VERTICES =================
\node[vertex] at (-0.6,-0.3) {};
\node[vertex] at (0.6,-0.2) {};
\node[vertex] at (0,0.7) {};

\node at (0,-1.8) {\small Edgeless};

% ================= PATH =================
\node[vertex] (p1) at (\dx-1.2,0) {};
\node[vertex] (p2) at (\dx-0.4,0.2) {};
\node[vertex] (p3) at (\dx+0.4,-0.1) {};
\node[vertex] (p4) at (\dx+1.2,0.1) {};

\draw[edge] (p1) -- (p2);
\draw[edge] (p2) -- (p3);
\draw[edge] (p3) -- (p4);

\node at (\dx,-1.8) {\small Path ($P_4$)};

% ================= K5 =================
\begin{scope}[shift={(2*\dx,0)}]

\node[vertex] (k1) at (90:1.3) {};
\node[vertex] (k2) at (18:1.3) {};
\node[vertex] (k3) at (-54:1.3) {};
\node[vertex] (k4) at (-126:1.3) {};
\node[vertex] (k5) at (162:1.3) {};

\draw[edge] (k1)--(k2) (k1)--(k3) (k1)--(k4) (k1)--(k5)
            (k2)--(k3) (k2)--(k4) (k2)--(k5)
            (k3)--(k4) (k3)--(k5)
            (k4)--(k5);

\end{scope}

\node at (2*\dx,-1.8) {\small Complete ($K_5$)};

% ================= MULTIGRAPH =================
\node[vertex] (m1) at (3*\dx-0.8,0) {};
\node[vertex] (m2) at (3*\dx+0.8,0) {};

\draw[edge] (m1) to[bend left=25] (m2);
\draw[edge] (m1) to[bend right=25] (m2);
\draw[edge] (m1) -- (m2);

\node at (3*\dx,-1.8) {\small Multigraph};

\end{tikzpicture}
    \caption{Three examples of simple graphs and one example of a multi-graph.}
    \label{fig:graphs}
\end{figure}
The maps Guthrie discussed correspond to a special type of graphs called planar graphs. A \textit{planar graph} is a graph that can be drawn without edge-crossings, i.e. it can be embedded in the 2-dimensional plane. Every graph shown in Figure~\ref{fig:graphs} is planar except the complete graph $K_5.$ Thus, Guthrie's observation can be stated more formally:
\begin{theorem}[The Four Color Theorem]\label{four_color_thm}
    For any planar graph $G,$ the chromatic number $\chi(G)$ of $G$ is at most $4$.
\end{theorem}
Here, the \textit{chromatic number} of a graph $G$, denoted $\chi(G)$, is the smallest number of colors required to color the vertices of a graph such that no two adjacent vertices share a color. Vertices are considered \textit{adjacent} if there is an edge between them.

Though the first written statement of Theorem~\ref{four_color_thm} appeared in 1852, it went largely unnoticed beyond a few correspondences until mathematician Arthur Cayley published a short article regarding the problem in 1878 as mentioned in~\cite{cayley}. Before the end of the decade, a man named Sir Alfred Kempe had published a ``complete'' proof~\cite{FCHistory}. However, 11 years after this supposed proof had been accepted by the mathematical community, it was disproved by Percy Heawood~\cite{heawood1898}. While unable to rectify Kempe's proof, he readily extracted a proof of the five color theorem. Additionally, Kempe's novel chain construction would become central to edge colorings. Perhaps the two most famous uses of these chains were D\'enes K\H{o}nig's and Vadim Vizing's publications bounding the chromatic indices of finite graphs \cite{konig, VizingProof} as we will mention in Section~\ref{section motivation}. Meanwhile, allured by its child-like simplicity, mathematicians would spend more than a century researching the four color theorem before the first viable proofs emerged with the assistance of computers. 

In 1976, Appel and Haken invested over 1000 hours of computer time to check 1,936 configurations and finally verified the four color theorem~\cite{Appel4Color}. However, their proof was met with controversy and hesitation: this was the first major mathematical result to be solved with the help of a computer. Mathematicians struggled to accept a proof which could not be personally checked by peer reviewers and diverged so notably from the mainstream mathematical reasoning found in proofs~\cite{suffices}. Following its publication, the mathematical community rapidly began sharing revisions, each requiring fewer computer-assisted cases, desperate to solidify the result.

This solution and its subsequent simplifications, with their reliance on hand-written computer code subject to programmer error and a history scattered with false claims, demonstrate both the beauty and the necessity of formal verification. With the development of interactive theorem provers, also known as proof assistants, one needs only to rely on the system itself, and often times even the kernels of these systems have themselves been formally verified\footnote{See, e.g., the MetaRocq project~\cite{meta}.}. In addition, the necessary detail required to formalize the proof alleviated some of the disappointment with Appel and Haken's original proof, that computers buried mathematical insight behind ones and zeros. As Georges Gonthier remarked, ``to produce a formal proof one must make explicit every single logical step of a proof; this both provides new insight in the structure of the proof, and forces one to use this insight to discover every possible symmetry, simplification, and generalization,''~\cite{Gonthier}. That is, the author must intimately understand the proof to formalize it. In 2008, Gonthier published a fully verified proof of the four color theorem~\cite{Gonthier} as the first substantial mathematical result to be implemented in a proof assistant. Today, the four color theorem is widely accepted and Gonthier's proof is in a well-maintained open-source library. Thus, while the impacts of this theorem continue to shape mathematics today, we conclude our short tale by encouraging the reader to glance at Gonthier's formalization in Rocq~\cite{4CRocq}. 

\section{Motivation}\label{section motivation}
As mathematicians labored over the four color theorem, a plethora of partial results, adjacent theorems, and new proof techniques established graph coloring as its own subfield. In addition to vertex colorings, there is also substantial research surrounding edge colorings, total colorings, balanced colorings, and fractional colorings, to name a few. This thesis will focus specifically on edge colorings. Similar to a vertex coloring, an \textit{edge coloring} assigns colors to the edges of a graph. In fact, edge colorings can be regarded as special cases of vertex colorings. The \textit{line graph} of a graph $G$, also commonly referred to as the conjugate or edge-to-vertex dual, can be constructed by associating each edge with a vertex and connecting the vertices if and only if the corresponding edges share a common vertex in $G.$ Then an edge coloring of a graph is exactly the vertex coloring of its line graph. 

The \textit{chromatic index} of a graph $G,$ denoted $\ci(G)$ or simply $\ci$ when $G$ is inferred and defined analogously to the chromatic number, is the smallest number of colors required to color the edges of a graph such that no two adjacent edges share a color. Edges are considered \textit{adjacent} if they share a common vertex. In 1964 and 1967, respectively, Vadim Vizing\cite{VizingProof} and Gupta\cite{Gupta} independently published the most famous theorem in edge coloring:
\begin{theorem}[Vizing's Theorem]\label{Vizing} 
    For any finite, simple graph $G$, we have $\Delta(G) \leq \ci(G) \leq \Delta(G)+1$
\end{theorem}
Here, the maximum degree $\Delta(G)$ or $\Delta$ of a graph $G$ is the maximum number of edges at any vertex in $V(G)$. Vizing's theorem imposes a much stricter bound on the chromatic index than exists for the chromatic number of an arbitrary graph. We saw that at most four colors were required in the vertex coloring of a planar graph; for non-planar graphs, Brooks showed $\,\chi\,$ is between $2$ and $\Delta + 1$ \cite{chromatic_number}. With such a range, it is all the more shocking the chromatic index is restricted to only two values. This allows us to classify graphs as Class-I if $\ci = \Delta$ and Class-II if $\ci = \Delta+1.$  K\H{o}nig obtained the first major edge coloring result in 1916 when he proved \textit{bipartite} graphs \cite{konig}, that is, graphs with chromatic number $2,$ are always Class-I. However, despite the apparent simplicity in contrast to vertex colorings, there are many open problems on characterizing Class-II graphs~\cite{class2}, and classifying a graph as Class-I or Class-II is NP-hard for arbitrary graphs~\cite{Ian}.

According to statistics collected by Crescenzi and Kann, graph coloring is a central issue in combinatorial research on NP-complete and NP-hard problems today \cite{NPStats}. Furthermore, graph coloring has numerous applications in both formal and applied mathematics; for a more comprehensive list, see the survey by Thakare and Bhapkar \cite{Survey} containing popular results and open problems related to graph coloring as well as references to some of the field's applications. Edge coloring in particular appears in route planning, transfer assignment in computer networks, latin squares, matrix algebra, and quantum field theory, to name a few \cite{random_app, LatinSquareCameron, LatinSquareHilton, chinesepost, bigSurvey}.

Meanwhile, formal verification systems have a wide range of applications from pure mathematics to systems engineering and artificial intelligence. Notable examples include the odd-order theorem ~\cite{groupTheory}, the verified C compiler CompCert\cite{compCert}, and recent efforts to integrate formal reasoning into artificial intelligence systems at companies like Axiom~\cite{axiom} and Harmonic~\cite{harmonic}. Gonthier's work demonstrates the potential of integrating these proof assistants with graph theory. In fact, most of these systems today already maintain graph theory libraries. For instance, Lars Noschinski wrote a library for directed graphs in Isabelle~\cite{graphLibIzzy}, the Lean Mathematical Library contains a graph theory component~\cite{mathlib}, and Rocq itself has two well-developed theoretical graph libraries~\cite{mathcomp, graphLibRock}. Edge coloring is a significant field of research, but the current scope of Rocq libraries do not yet include it. Our \fxn{edge-coloring} file defines an edge coloring function as well as basic actions and lemmas regarding recoloring and swapping colors, absent sets, and trivial colorings. Integrating this with the library would greatly benefit the Rocq community by establishing a standard definition capable of utilizing the rich set of graph results already present. Adding Kempe chains and a proof of Vizing's theorem, from which we can conclude all graphs are either Class-I or Class-II, to the library would also support many future endeavors including the open problems we alluded to above.
% This ongoing work intends to reduce the gap between what is already informally proved and what is not, such that other graph theorists may have a framework to formalize their results, especially when their proofs require the analysis of dozens of mechanical cases (the Four-Color Theorem is an example of a result involving an overwhelming number of cases [7]).

% on the one hand, to get confident about the proof and, on the other, to have the advantage that a reader can accept it without the need of manually checking step by step (it eventually could reduce the time spent in the peer-review process

\section{Related Work}\label{section related}
Both Vizing's and Gupta's proofs, while correct, lacked simplicity. As such, several variations on the original works have been published in an attempt to find a more elegant approach. Bollob\'as sketched a proof in 1979~\cite{boll}, Rao and Dijkstra developed algorithms with a proof of termination and auxiliary data structures~\cite{dijk}, and Misra and Gries further simplified these~\cite{Misra}. Many expository works and introductory graph theory text books have also compiled a proof. In particular, we were inspired by \cite{Green}, \cite{Diestel}, and \cite{GECBook}. Even within these publications, the proof that $\ci$ is $\Delta$ or $\Delta + 1$ remains rather involved, often dependent on fragile relationships between specialized data structures. Additionally, proofs gloss over some of the finer-grained details and omit non-trivial cases, intentionally or unintentionally. This, coupled with the wide literature on edge colorings, prompted the formalization of Vizing's theorem in first Rodin, then Lean, and now Rocq. In the same way variations on Vizing's original proof hold value even though the result is known, it is beneficial to formalize it multiple times for a variety of reasons. Firstly, it strengthens the implementation when it spans multiple languages and is not reliant on a single system. For the remainder of this section, we will discuss our project in comparison to related works.

In 2011, Joachim Breitner modeled Misra and Gries's algorithm in Event-B with verification in Rodin\cite{Breitner}. While sound according to the system, the individual lemmas in Rodin are write-once, read-never (WORN) and, most importantly, they rely on the soundness of Rodin itself. Although this is the case for every implementation into a formalization system, reliance on the system's soundness is especially bold with Rodin because it has a history of errors. Breitner himself acknowledged ``relying only on the Rodin platform for verifying the proof obligations would not constitute a rigorous proof.''

Undergraduates Aarbi, He, and Gatti wrote a structural template to formalize Vizing's theorem in Lean4 with Mathlib as part of a spring 2025 research project at the University of Washington~\cite{math480}. Before Bhoja's publication in December of 2025, this was our only reference for an attempted formalization of Vizing's theorem in a language based on the calculus of inductive constructions. Their definition of edge colorings and use of a dependent color type supported our final decision to represent edge colorings similarly. Additionally, we appreciated their inductive set-up and the creation of an explicit coloring to trivially solve the base case. Of course, their project merely outlined a potential approach and consisted mostly of comments and statements with empty proofs.

The original vision for this thesis was to construct the first complete proof of Vizing's theorem in a verified proof assistant. A little over half-way through our project, Bhoja published a complete Lean implementation of Misra and Grieg's algorithm in December of 2025\cite{Bhoja}. While this thesis shares similarities to Bhoja's approach, there are several key differences. Most importantly, we wish to offer formalized edge colorings to the Rocq community. Bhoja sought to create an efficient polynomial-time algorithm and as such did not work intending to contribute to Lean's existing graph theory libraries, instead creating specialized data structures from scratch. We, on the other hand, though valuing efficiency in our design, acted with the goal of extending Rocq's libraries as discussed in Section~\ref{section motivation}. Also, migrating the mathematical proofs into Rocq was non-trivial, even with the example in Lean, for, while Lean and Rocq have similar foundational logic, they differ in syntax as well as current results and definitions within their respective libraries, including the encouragement of entirely different graph structures. We also provide greater detail in our written proofs to more closely align with the code in Rocq, and in Chapter~\ref{chapter implement} we elaborate on implementation decisions made, especially in comparison to Bhoja's method. Finally, we formalize the lower bound of Vizing's theorem which, though shorter, is not included in any of the other projects we have mentioned.

\section{Background and Outline}
As mentioned, a graph $G = (V, E\,)$ refers to a set of vertices $V(G)$ and a set of edges $E(G)$ between these vertices. For any two vertices $x, y$ in $V(G),$ we will use the pair notation $(x, y)$ to denote an edge connecting $x$ and $y.$ All graphs hereafter are considered finite, simple, and undirected unless otherwise specified. That is to say, the set $V(G)$ is finite, there is at most one edge between any two vertices, there are no self-loops, and for any two vertices $x, y$ in $V(G)$ such that $\e{x}{y}$ is in $E(G),$ we have $\e{y}{x}$ in $E(G).$

This thesis examines the edge coloring of simple graphs with formal verification. In particular, we prove Vizing's theorem on the chromatic index of a graph in the proof assistant Rocq. \textit{Formal verification} is the application of formalized mathematical methods to prove the correctness of a statement or system. \textit{Proof assistants} are software systems that employ formal methods to allow us to construct proofs as computer programs and mechanically verify their correctness. The validity of a proof is migrated from humans to the correctness of the proof assistant itself. They are especially powerful for proofs involving many cases, as seen with the four color theorem. The guarantees offered by proof assistants can be a strong foundation for applications that require a high level of certainty such as security systems, computer kernels, and hardware design. Examples of these systems include Rocq, Lean, Isabelle, Agda, Mizar, and Metamath. First released in 1989, Rocq (formerly Coq) has gained traction in the computer science community. It is constructive in nature (that is, it does not assume the axiom of the excluded middle), has well-developed libraries for large-scale mathematical proofs such as Gonthier's four color theorem, and supports the immediate extraction of verified algorithms to runnable programs. 

In the remainder of this thesis, we will describe our implementation of Vizing's theorem in Rocq. In Chapter~\ref{chapter math}, we present our own complete proof of Vizing's theorem and supporting results. As much as is reasonable, we try to align these proofs with those in Rocq. We also offer commentary relating these definitions and lemmas to the implementation process, though Chapter~\ref{chapter implement} explores the details and challenges of our execution in greater depth. Specifically, Chapter~\ref{chapter implement} focuses on the project's main contributions and the most architecturally significant sections in the code. Finally, we summarize the project in the Conclusion, touching on potential directions for future work. The Rocq code is displayed in Appendix~\ref{AppendixB} and can also be accessed from our GitHub repository~\cite{my_code}.
